% Appendix E (writer: steep-a). Proofs for sec:split. Source: steeper notes-final §3.1-§3.3 (Lemmas 3.2, 3.2', Prop 3.3, Thm 4).
\section{Proofs for the model--randomness split}\label{app:split}

The proof of \cref{thm:split:lambda1} is given in the main text.

\subsection{Proof of the Ville lemma}\label{app:split:ville}

We prove \cref{lem:split:ville}. Let $Z_n:=\xi(D_n)/\nu(D_n)$. Given the past and the next input $x_{n+1}$, which the input rule chooses before $y_{n+1}$ is drawn,
\[
\Ex[Z_{n+1}]=\sum_y\nu(y\mid x_{n+1})\frac{\xi(D_n+(x_{n+1},y))}{\nu(D_n)\,\nu(y\mid x_{n+1})}\le\frac{\sum_y\xi(D_n+(x_{n+1},y))}{\nu(D_n)}\le Z_n,
\]
the sum running over labels of positive $\nu$-probability, by the one-point inequality. So $(Z_n)$ is a nonnegative supermartingale with $Z_0=\xi(\emptyset)\le1$, also when the input rule is adaptive or randomised. Ville's maximal inequality \citep{ville1939etude} gives $\Prob[\sup_nZ_n\ge2^j]\le Z_0\,2^{-j}\le2^{-j}$. For the consequence, $\xi_{\mathcal C}\ge2^{-\ell(\pi)}\nu_\pi$, so $\nu_\pi(D_n)\le2^{\ell(\pi)}\xi_{\mathcal C}(D_n)\le2^{\ell(\pi)+\delta^*}\nu(D_n)$. \qed

\subsection{Proof of the stream-prefix lemma}\label{app:split:stream}

We prove \cref{lem:split:stream}. Order $Y=\{y^{(1)},\dots,y^{(|Y|)}\}$.

\emph{Private streams.} Let $j(x)$ be the length-lexicographic index of $x$, so $j(x)\le2^{d+1}-2$ for $|x|\le d$, and let $\langle a,b\rangle:=(a+b)(a+b+1)/2+b$ be the Cantor pairing, which is injective and increasing in $b$. For a finite string $w$ (bits indexed from $0$), the \emph{stream prefix} of $x$ in $w$ is $w_{\langle j(x),0\rangle}w_{\langle j(x),1\rangle}\cdots$, continued while the positions are below $|w|$. Distinct inputs read disjoint positions.

\emph{Certification.} For an input $x$ put $F_h:=\sum_{h'\le h}\nu(y^{(h')}\mid x)$, so $F_0=0$ and $F_{|Y|}=1$. At depth $i\ge1$ compute, by a fixed algorithm from the program for $\nu$, rationals $q_{i,h}$ with $|q_{i,h}-F_h|\le2^{-i-2}$ ($q_{i,0}=0$, $q_{i,|Y|}=1$). A stream prefix $b_1\cdots b_i$ names the dyadic interval $I=[0.b_1\cdots b_i,\ 0.b_1\cdots b_i+2^{-i})$. It \emph{certifies} $y^{(h)}$ at depth $i$ if $q_{i,h-1}+2^{-i-2}\le\min I$ and $\sup I\le q_{i,h}-2^{-i-2}$. Then $I\subseteq[F_{h-1},F_h]$.

\emph{Consistency.} Suppose the stream of $x$ certifies $y^{(h)}$ at depth $i$ and $y^{(h')}$ at depth $i'$, with $h<h'$. The two intervals are nested, since one stream prefix extends the other; the inner one has positive length and lies in $[F_{h-1},F_h]\cap[F_{h'-1},F_{h'}]$. For $h<h'$ this intersection is empty or the single point $F_h$, a contradiction. So $\varphi_w(x)$, defined as the label certified at some depth within the stream prefix of $x$ in $w$ (undefined if none), is a well-defined partial computable function, and an index of it is computable from $(V(\pi_\nu),w)$.

\emph{Probability.} Every depth-$i$ dyadic interval inside $[F_{h-1}+2^{-i-1},F_h-2^{-i-1}]$ certifies $y^{(h)}$, because $q_{i,h-1}+2^{-i-2}\le F_{h-1}+2^{-i-1}$ and $q_{i,h}-2^{-i-2}\ge F_h-2^{-i-1}$. These intervals cover that segment except for at most $2\cdot2^{-i}$, and the segment has length $\nu(y^{(h)}\mid x)-2^{-i}$. So a uniform stream prefix of length at least $i$ certifies $y^{(h)}$, and hence gives $\varphi_w(x)=y^{(h)}$, with probability at least $\nu(y^{(h)}\mid x)-3\cdot2^{-i}$.

\emph{Independence.} For $w$ uniform on $\{0,1\}^L$, the stream prefixes of distinct inputs are independent and uniform. If $L\ge(2^{d+1}+i)^2$, every input of length at most $d$ has a stream prefix of length at least $i$, because $\langle j(x),i-1\rangle<(j(x)+i)^2\le(2^{d+1}+i)^2$. Take $i=i_D$; then $3\cdot2^{-i_D}\le\beta_D/(2n)$, and
\[
\Pr_w[\varphi_w\supseteq D]\ \ge\ \prod_{(x,y)\in D}\big(\nu(y\mid x)-3\cdot2^{-i_D}\big)\ \ge\ \nu(D)\Big(1-\frac1{2n}\Big)^n\ \ge\ \frac{\nu(D)}2,
\]
using $\nu(y\mid x)\ge\beta_D$ and Bernoulli's inequality.

\emph{Programs.} Let $V'$ be the prefix machine that reads $\pi_\nu$ (by simulating $V$, which consumes exactly $\pi_\nu$), then a self-delimiting code of an integer $t$ of length at most $2\log(t+1)+1$, then $L_t:=2^{2^t}$ bits $w$, and outputs an index of the deterministic conditional semimeasure $\varphi_w$. By closure under prefix-free simulation (\cref{def:setup:class}) there is $r$ with $V(ru)=V'(u)$. The programs $p_w:=r\,\pi_\nu\,\code(t)\,w$ are distinct for distinct $w$. With $t=t_D$ we have $2^{t_D}\ge2d+2+2\log(i_D+1)$, so $L_{t_D}\ge\big(2^{d+1}(i_D+1)\big)^2\ge(2^{d+1}+i_D)^2$.

\emph{Sum.}
\[
\Wdet(D)\ge\sum_w2^{-|p_w|}\ind[\varphi_w\supseteq D]=2^{-|r|-k-|\code(t_D)|}\Pr_w[\varphi_w\supseteq D]\ge2^{-k-2\log(t_D+1)-|r|-2}\,\nu(D).
\]
\qed

\emph{Computation.} A direct computation (code in the accompanying repository, seed 7) checked the certification step on 600 random samplers ($|Y|\le4$, rational values with various denominators, Bern$(1/3)$ included) with adversarial approximations $q_{i,h}=F_h\pm2^{-i-2}$: in all 1646 label--depth cases the certified mass was at least $\nu-3\cdot2^{-i}$, no two certificates conflicted, and the end-to-end ratio at depth $i_D$ was at least $1/2$ (smallest $0.645$). It also checked the pairing and length facts for $d\le12$, $i\le64$. The cost $2\log(t_D+1)$ is $5.6$ bits at $d=20$, $n=10^3$ ($\beta_D=1/3$); $6.3$ bits at $d=100$, $n=10^6$ (for $\beta_D=1/3$ and for $\beta_D=2^{-10}$); $7.2$ bits at $d=1000$, $n=10^{12}$ ($\beta_D=10^{-3}$); and $8.9$ bits at $d=10^6$, $n=10^{15}$ ($\beta_D=1/3$). The hashing cost of \cref{lem:split:hash} at the same sizes is $31.9$, $56.3$ (with $r=2$) or $60.1$ (with $r=10$), $106.6$ and $142.7$ bits.

\emph{Ville's inequality in simulation.} For \cref{lem:split:ville}, a simulation with the true sampler (code in the accompanying repository, seed 4; 4000 runs) gave empirical frequencies $\Prob[\delta^*\ge j]=0.119,0.055,0.032,0.014,0.003$ for $j=1,2,3,4,6$, and an independent simulation with an adversarial adaptive input rule, which puts the next input where $\xi$ deviates most from $\nu$, gave $0.355,0.163,0.069,0.035,0.0065$; all are below $2^{-j}$.

\subsection{An alternative for dyadic samplers: hashing}\label{app:split:hash}

\begin{lemma}[Derandomising a dyadic sampler by hashing]\status{proved; computed}\src{steeper Lemma 3.2; checks c4(B), c6(E); referee r5, q3(A)}\label{lem:split:hash}
Let the class be $\Hcl{\times}$, let $\nu$ be a sampler whose values are multiples of $2^{-r}$, with a program $\pi_\nu$ of length $k$, and let $D$ be a data set with $n$ inputs, all of length at most $d$. Then
\[
\Wdet(D)\ \ge\ 2^{-k-2\log(n+1)-2\log(d+1)-2\log(r+1)-\kappa}\,\nu(D).
\]
\end{lemma}

\begin{proof}
\emph{Exact values.} From the index $V(\pi_\nu)$ and $r$ one computes $\nu$ exactly: enumerate lower approximations $\nu_s(y\mid x)\uparrow\nu(y\mid x)$ until $\sum_y\nu_s(y\mid x)>1-2^{-r}$, which happens because $\nu$ is a measure and $Y$ is finite. Then every deficit $\nu-\nu_s$ is below $2^{-r}$, and $\nu(y\mid x)$ is the least multiple of $2^{-r}$ that is at least $\nu_s(y\mid x)$.

\emph{Canonical sampler.} For $b\in\{0,\dots,2^r-1\}$ let $\sigma(x,b)$ be the label $y$ whose cumulative interval $[\sum_{y'<y}\nu(y'\mid x),\sum_{y'\le y}\nu(y'\mid x))$ contains $b2^{-r}$. Then $\Pr_b[\sigma(x,b)=y]=\nu(y\mid x)$ for uniform $b$.

\emph{Hash.} Let $q:=\max(d+1,r)$, fix (computably from $q$) an irreducible polynomial of degree $q$ over GF$(2)$, and embed $\{0,1\}^{\le d}$ injectively into GF$(2^q)$ by mapping $x$ to the element with coefficient vector $1x$. For a seed $s\in\mathrm{GF}(2^q)^n$ let $h_s(x)$ be the first $r$ coefficient bits of $\sum_{j<n}s_j\,x^j$, computed in GF$(2^q)$ with $x$ embedded. For distinct $x_1,\dots,x_n$ the map $s\mapsto(\sum_js_jx_i^j)_{i\le n}$ is a bijection (a Vandermonde matrix with distinct nodes), so for uniform $s$ the values $h_s(x_i)$ are independent and uniform on $\{0,1\}^r$.

\emph{Programs, one per seed.} Let $V'$ be the prefix machine that reads $\pi_\nu$ (by simulating $V$), then self-delimiting codes of $n$, $d$ and $r$, then $qn$ bits $s$, and outputs an index of the deterministic conditional semimeasure $x\mapsto\sigma(x,h_s(x))$ on $\{0,1\}^{\le d}$, undefined elsewhere. By closure under prefix-free simulation there is $r'$ with $V(r'u)=V'(u)$. The programs $p_s:=r'\pi_\nu\code(n)\code(d)\code(r)\,s$ are distinct for distinct $s$, deterministic, and of length $L_0+qn$ with $L_0\le k+2\log(n+1)+2\log(d+1)+2\log(r+1)+\kappa$.

\emph{Sum.} $\sum_s2^{-|p_s|}\ind[\nu_{p_s}\supseteq D]=2^{-L_0}\Pr_s[\forall i:\sigma(x_i,h_s(x_i))=y_i]=2^{-L_0}\prod_i\nu(y_i\mid x_i)=2^{-L_0}\nu(D)$.
\end{proof}

Distinct seeds often give the same function on $X$; the sum is over programs, one per seed, so this does not matter. A direct computation (code in the accompanying repository) checked the independence exactly over GF$(16)$ for $n\le3$ (3232 cases) and the full statement with the embedding $1x$, $q=\max(d+1,r)$, $r\in\{2,3,4\}$, $|Y|=3$ and $n\le3$ (2106 cases), with maximum error $0$. The hashing idea for fair coins and the program mixture is in \citet[the hashing remark]{claudeSFI2026fundamentals}.

\subsection{Proof of the SI proposition}\label{app:split:si}

We prove \cref{prop:split:si}. (i) Under $\nu$, $Z_n:=M(z_{1:n})/\nu(z_{1:n})$ is a nonnegative supermartingale with $Z_0\le1$, because $M(z0)+M(z1)\le M(z)$; Ville's inequality gives the tail bound, and $\Pi_1(\pi_\nu\mid z_{1:n})=2^{-k}\nu(z_{1:n})/M(z_{1:n})\ge2^{-k-\delta^*}$. (ii) Let $w:=z_{1:n}0^\infty$ and let $\delta_w$ be the point mass on its prefixes. Given $n$ and $\pi_\nu$, the restriction of $\nu$ to $\{0,1\}^n$ is a computable measure, so (F2) gives $K(z_{1:n}\mid n,\pi_\nu)\le-\log\nu(z_{1:n})+\kappa$, and (F1), (F3) give $K(z_{1:n})\le K(n)+k-\log\nu(z_{1:n})+\kappa$. An index of $\delta_w$ is computable from $z_{1:n}$, so the class contains a program for $\delta_w$ of length at most $K(z_{1:n})+\kappa$, and $\delta_w$ is consistent with $z_{1:n}$. (iii) The programs ``a fixed decoder, a shortest program for $\nu$, then a Shannon--Fano code of $z$ under $\nu$'' have total weight at least $2^{-k-\kappa}\nu(z)$; this is the usual proof of $M\ge2^{-K(\nu)-O(1)}\nu$ \citep{li2008introduction}. \qed

\subsection{Proof of the private-randomness theorem}\label{app:split:private}

We prove \cref{thm:split:private}. Write $D=D_n$.

\emph{(i).} $J_\lambda(\hat\pi;D)\le J_\lambda(\pi_\nu;D)=\lambda k-\log\nu(D)$. By \cref{lem:split:ville}, $-\log\nu_{\hat\pi}(D)\ge-\log\nu(D)-\ell(\hat\pi)-\delta^*$. Substituting gives $(\lambda-1)\ell(\hat\pi)\le\lambda k+\delta^*$. The upper bound on the loss is $J_\lambda(\hat\pi;D)\le J_\lambda(\pi_\nu;D)$ with $\lambda\ell(\hat\pi)\ge0$.

\emph{(ii).} $\sum_{\ell(\pi)>L}2^{-\lambda\ell(\pi)}\nu_\pi(D)\le2^{-(\lambda-1)L}\xi_{\mathcal C}(D)\le2^{-(\lambda-1)L+\delta^*}\nu(D)$, and the normaliser $\xi_\lambda(D)$ is at least the sampler's term $2^{-\lambda k}\nu(D)$.

\emph{(iii).} A consistent deterministic $\pi$ has $1=\nu_\pi(D)\le2^{\ell(\pi)+\delta^*}\nu(D)$, which is the length bound. For the power posterior, $\sum_{\det}2^{-\lambda\ell(\pi)}=\sum_{\det}2^{-\ell(\pi)}2^{-(\lambda-1)\ell(\pi)}\le2^{-(\lambda-1)(-\log\nu(D)-\delta^*)}\Wdet(D)$, the sums running over consistent deterministic programs; with $\Wdet\le\xi_{\mathcal C}\le2^{\delta^*}\nu(D)$ and the normaliser at least $2^{-\lambda k}\nu(D)$ this gives the bound. For the tempered posterior, the fill-ins' total weight is $\Wdet(D)\le2^{\delta^*}\nu(D)$, and the normaliser is at least the sampler's term $2^{-k}\nu(D)^\eta$.

\emph{(iv), likelihood ratios.} Fix $\pi\in\mathcal C$ and let $\Lambda_n:=\nu_\pi(D_n)/\nu(D_n)$ and $g_n:=\Hel(\nu(\cdot\mid x_n),\nu_\pi(\cdot\mid x_n))$; $g_n$ is determined before $y_n$ is drawn. Given the past and $x_{n+1}$,
\[
\Ex\sqrt{\Lambda_{n+1}/\Lambda_n}=\sum_y\nu(y\mid x_{n+1})\sqrt{\nu_\pi(y\mid x_{n+1})/\nu(y\mid x_{n+1})}=1-g_{n+1}.
\]
If $g_{n+1}=1$, then almost surely $\nu_\pi(y_{n+1}\mid x_{n+1})=0$ and $\Lambda$ stays $0$ from then on. Let $M_n:=\sqrt{\Lambda_n}/\prod_{i\le n}(1-g_i)$ while the product is positive, and $M_n:=0$ afterwards. Then $(M_n)$ is a nonnegative supermartingale, so it converges almost surely to a finite limit. On the event $\{\sum_ig_i=\infty\}$ either some $g_i=1$, and $\Lambda_n=0$ eventually, or $\prod_{i\le n}(1-g_i)\to0$, and $\sqrt{\Lambda_n}=M_n\prod_{i\le n}(1-g_i)\to0$. This is the argument behind Kakutani's dichotomy \citep{kakutani1948equivalence}; it holds for adaptive inputs because $g_n$ is predictable.

\emph{(iv), MAPs.} There are countably many programs, so almost surely, simultaneously, $\delta^*<\infty$ and $\Lambda_n(\pi)\to0$ for every $\pi$ that is not $\nu$-equivalent on the run. By (i), every $\lambda$-MAP lies in the finite set $\{\pi:\ell(\pi)\le L^*\}$ with $L^*:=(\lambda k+\delta^*)/(\lambda-1)$. A $\lambda$-MAP $\hat\pi$ has $\log\Lambda_n(\hat\pi)\ge\lambda\ell(\hat\pi)-\lambda k\ge-\lambda k$, from $J_\lambda(\hat\pi)\le J_\lambda(\pi_\nu)$. Each non-equivalent program in the finite set eventually has $\Lambda_n<2^{-\lambda k}$, so eventually none of them is a $\lambda$-MAP.

\emph{(iv), posterior.} For every $L$, using (ii) for the long programs and the sampler's term for the normaliser,
\[
\Pi_\lambda(\text{not equivalent}\mid D_n)\ \le\ 2^{\lambda k}\sum_{\ell(\pi)\le L,\ \text{not equivalent}}2^{-\lambda\ell(\pi)}\Lambda_n(\pi)\ +\ 2^{\lambda k+\delta^*-(\lambda-1)L}.
\]
The first term is a finite sum of terms that tend to $0$. Let $n\to\infty$, then $L\to\infty$. \qed

\emph{Computation.} In an explicit Kraft-valid product class (Kraft sum $0.909$) of samplers, fill-ins and partial memorisers, an independent computation (code in the accompanying repository, seed 777) ran 30 runs at 12 sample sizes and 4 values of $\lambda$ and found no violation of (i)--(iii) or of the upper end of \cref{thm:split:lambda1}(iii) in 1440 MAP checks.

\subsection{Proof of the shared-randomness theorem}\label{app:split:shared}

We prove \cref{thm:split:shared}. (i) By (Dom), $-\log U(o)\le k+\kappa_0-\log P_{\pi_\nu}(o)$, so $J_\lambda(\pi_U;o)=\lambda\ell_U-\log U(o)\le\lambda\ell_U+\kappa_0+k-\log P_{\pi_\nu}(o)$, which is below $\lambda k-\log P_{\pi_\nu}(o)=J_\lambda(\pi_\nu;o)$ exactly when $(\lambda-1)k>\lambda\ell_U+\kappa_0$; the same condition says $\ell^*_\lambda<k$, and every $\lambda$-MAP has length at most $\ell^*_\lambda$ by \cref{thm:absorb:absorption}(a). (ii) $\sum_{\pi\in A}2^{-\lambda\ell(\pi)}P_\pi(o)\le2^{-(\lambda-1)\min_A\ell}\sum_\pi2^{-\ell(\pi)}P_\pi(o)=2^{-(\lambda-1)\min_A\ell}\,\xi_1(o)\le2^{-(\lambda-1)\min_A\ell+\kappa_1}U(o)$, and the normaliser $\xi_\lambda(o)$ is at least $2^{-\lambda\ell_U}U(o)$. For SI with adaptive or non-computable inputs, Solomonoff's $M$ is a semimeasure on the whole interleaved string and need not dominate laws whose inputs are not modelled; $\MH$ dominates every lsc chronological semimeasure on every input sequence \citep{claudeSFI2026fundamentals}, so (Dom) holds in $\Hcl{\mathrm{chr}}$. \qed

In an explicit shared class that contains its own mixture, an independent computation (code in the accompanying repository, seed 4242) found no violation of (ii) or of \cref{thm:absorb:absorption}(a) and (c1) in 300 MAP checks.
